\documentclass[10pt]{article}
\usepackage[margin=2cm]{geometry}
\usepackage{booktabs}
\usepackage{graphicx}
\usepackage{hyperref}
\usepackage{amsmath}

% Cells still waiting on a Delta run. Search for \tbd to find them all.
\newcommand{\tbd}{\textbf{--}}

\begin{document}

\title{CSCI 5942: AI Engineering Assignment 4\\Distributed Training}
\author{Bill Hsu}
\maketitle

% =====================================================================
\section*{Setup}

All runs used NCSA Delta \texttt{gpuA40x4} nodes (four NVIDIA A40s per
node, 44\,GiB each), the \texttt{pytorch-conda/2.12} module, and bf16
autocast. Parameters, gradients and AdamW moments are fp32. Throughput is
the global batch divided by the median steady-state step time, as reported
by \texttt{trainer/train.py}. Correctness was checked on CPU with the
\texttt{gloo} backend: \texttt{tests/test\_ddp.py} (10 passed),
\texttt{tests/test\_tp.py} (9 passed) and \texttt{tests/test\_pp.py}
(12 passed).

% =====================================================================
\section{Part 0: Communication Cost}

Table~\ref{tab:busbw} lists the all-reduce bus bandwidth measured by
\texttt{bench/allreduce.py} (fp32, median of 20 iterations). Small buffers
are latency-bound. Above about 128\,MiB the curve flattens. The
calculations below use the plateau values: \textbf{18.1\,GB/s at two A40s}
and \textbf{15.8\,GB/s at four}.

\begin{table}[htbp]
\caption{All-reduce bus bandwidth on one \texttt{gpuA40x4} node,
$\text{busbw} = \frac{2(P-1)}{P}\cdot\frac{N}{t}$.}
\label{tab:busbw}
\centering
\begin{tabular}{rrr}
\toprule
Buffer (MiB) & busbw, 2 A40s (GB/s) & busbw, 4 A40s (GB/s) \\
\midrule
1    & 8.27  & 7.68  \\
4    & 14.61 & 12.61 \\
16   & 17.09 & 15.13 \\
64   & 17.80 & 15.58 \\
256  & 18.14 & 15.75 \\
1024 & 18.18 & 15.81 \\
\bottomrule
\end{tabular}
\end{table}

% =====================================================================
\section{Part 1: Data Parallel}

\texttt{broadcast\_module\_state} broadcasts every parameter and buffer
from the first rank of the data-parallel group. \texttt{naive\_all\_reduce\_grads}
all-reduces each gradient after the backward pass and divides it by the
group size. \texttt{OverlappedDDP} registers a post-accumulate-grad hook on
every parameter. The hook divides the gradient by the group size and
launches an asynchronous all-reduce, so communication overlaps with the
rest of the backward pass. \texttt{finish\_gradient\_synchronization} waits
on every handle before the optimizer step. The hook divides
\emph{before} the all-reduce starts because \texttt{p.grad} must not be
modified while an asynchronous all-reduce is still using it. This gives
the same result, because $\sum_r g_r/P = (\sum_r g_r)/P$.

\begin{table}[htbp]
\caption{\texttt{medium} (406M parameters), global batch 128, on A40s.
Scaling efficiency is $\text{tok/s}_P / (P \cdot \text{tok/s}_1)$.}
\label{tab:ddp}
\centering
\begin{tabular}{rlrrr}
\toprule
GPUs & Strategy & Step time (s) & Tokens/s & Scaling eff. \\
\midrule
1 & ---        & 5.317 & 24{,}653 & 100\% \\
2 & naive      & 2.777 & 47{,}204 & 95.7\% \\
2 & overlapped & 2.752 & 47{,}625 & 96.6\% \\
4 & naive      & 1.516 & 86{,}483 & 87.7\% \\
4 & overlapped & 1.462 & 89{,}660 & 90.9\% \\
\bottomrule
\end{tabular}
\end{table}
% TODO(you): fill from out/medium-a40-{naive,overlapped}-x{1,2,4}/summary.json
% (median_step_seconds, tokens_per_second). At one GPU, naive and
% overlapped are the same code path (dp=1 skips the wrapper), so report it once.

\paragraph{Q1: Communication overhead of naive DP at four GPUs.}
Each step reduces one fp32 gradient per parameter:
$N = 406{,}161{,}408 \times 4\,\text{B} = 1.62\,\text{GB}$.
A ring all-reduce over $P$ ranks puts $\frac{2(P-1)}{P}N$ bytes on each
rank's link. At $P=4$, that is $1.5 \times 1.62 = 2.44$\,GB. Dividing by
the measured bus bandwidth gives the predicted communication time:
\[
t_\text{comm}(P{=}4) = \frac{1.5 \times 1.62\,\text{GB}}{15.8\,\text{GB/s}}
  \approx 0.154\,\text{s},
\qquad
t_\text{comm}(P{=}2) = \frac{1.0 \times 1.62\,\text{GB}}{18.1\,\text{GB/s}}
  \approx 0.090\,\text{s}.
\]
With perfect scaling, the four-GPU step would take
$t_1/4 = 5.317/4 = 1.329$\,s. The measured naive step took 1.516\,s, so the
overhead is 0.187\,s, compared with the predicted 0.154\,s. The prediction accounts
for 82\% of the measured overhead.
% TODO(you): fill in the blanks and say how close the prediction is.
% If the measured overhead is larger, name the rest: the all-reduce is
% issued as hundreds of per-parameter calls (small ones run at the
% latency-bound end of Table 1), plus the torch.cuda.synchronize after it.

\paragraph{Q2: How much does overlap hide?}
Overlapped saves $1.516 - 1.462 = 0.054$\,s per step relative to naive,
which is 35\% of the predicted 0.154\,s.
% TODO(you): the ceiling. Two points:
%  (1) Overlap can only hide communication behind backward compute that is
%      still running. The gradients of the first layers are ready only at
%      the very end of backward, so their all-reduce always runs after the
%      compute finishes.
%  (2) With gradient accumulation, only the last micro-batch's backward
%      triggers reductions (no_sync on the others). The work available to
%      hide behind is therefore one micro-batch's backward, not the whole step's.

% =====================================================================
\section{Part 2: Tensor Parallel}

$f$ is an identity in the forward pass and an all-reduce in the backward
pass. $\bar f$ is the mirror image. \texttt{ColumnParallelLinear} enters
through $f$ and multiplies the full input by its row-shard of $W$.
\texttt{RowParallelLinear} multiplies its input shard by its column-shard
of $W$, sums the partial results with $\bar f$, and only then adds the
bias. Adding the bias before the reduction would add it $P$ times. In
\texttt{TensorParallelMLP}, \texttt{c\_fc} is column-parallel and
\texttt{c\_proj} is row-parallel. GELU is elementwise, so it runs on the
shard, and each MLP needs one all-reduce in the forward pass and one in
the backward pass.

\begin{table}[htbp]
\caption{\texttt{large} (838M parameters), global batch 128, one
\texttt{gpuA40x4} node.}
\label{tab:tp}
\centering
\begin{tabular}{lrrrr}
\toprule
Configuration & Step time (s) & Tokens/s & MFU & Peak memory (GiB) \\
\midrule
1 GPU          & 11.283 & 11{,}617 & 40.4\% & 18.55 \\
DP=4           & 3.127  & 41{,}922 & 36.4\% & 18.55 \\
TP=4           & 12.100 & 10{,}833 & 9.4\%  & 6.34  \\
DP=2 $\times$ TP=2 & 6.036 & 21{,}715 & 18.9\% & 10.50 \\
\bottomrule
\end{tabular}
\end{table}
% TODO(you): fill from out/large-tp-{x1,dp4,tp4,dp2tp2}/summary.json, or
% from the "summaries" block at the end of logs/a4-tp-a40-<id>.out.

\paragraph{Q1: Ranking, and why TP=4 trails DP=4.}
% TODO(you): rank the four by tokens/s. Points to make:
%  - Same arithmetic, so the gap is communication that cannot be hidden.
%    TP's all-reduces sit on the critical path: the next layer needs the
%    reduced activation before it can start. DP's reduction overlaps with
%    backward.
%  - TP's benefit is memory. Each GPU holds 1/4 of the weights, gradients
%    and optimizer state. Point at the peak-memory column: TP=4 vs DP=4.
%  - Verdict: large fits on one A40 (see the 1 GPU row), so DP wins here.
%    TP is worth its cost only when the model does not fit.

\paragraph{Q2: Bytes on the wire per step.}
\emph{DP=4} all-reduces one fp32 gradient per parameter once per step:
$N_\text{DP} = 838{,}064{,}640 \times 4\,\text{B} = 3.35$\,GB.

\emph{TP=4} all-reduces activations. Each step processes
$128 \times 1024 = 131{,}072$ tokens of width $C = 1280$, so each
activation all-reduce, summed over all micro-batches, moves
$131{,}072 \times 1280 = 1.68\times10^{8}$ elements. Each of the 36
blocks has two TP regions (attention and MLP). Each region costs one
$\bar f$ all-reduce in the forward pass and one $f$ all-reduce in the
backward pass. The output head adds one more $f$ in the backward pass.
Under autocast, the forward payload is the bf16 output of a linear layer
(2\,B). The backward payload is the gradient of a LayerNorm output, which
is fp32 (4\,B). We checked both dtypes on a GPU.
\[
N_\text{TP} = 1.68\times10^{8}\times\bigl(72 \times 2\,\text{B} + 73 \times 4\,\text{B}\bigr)
  = 73.1\,\text{GB}.
\]
TP=4 therefore moves about $\mathbf{22\times}$ more data than DP=4. At
15.8\,GB/s, the predicted communication times are:
\[
t_\text{DP} = \frac{1.5 \times 3.35}{15.8} \approx 0.32\,\text{s},
\qquad
t_\text{TP} = \frac{1.5 \times 73.1}{15.8} \approx 6.9\,\text{s}.
\]
This is a lower bound for TP. With micro-batch 2, each TP all-reduce is
only 5--10\,MB. Table~\ref{tab:busbw} shows that buffers of this size get
about 13--14\,GB/s instead of 15.8\,GB/s, which raises the estimate to
about 7.6--8.4\,s.
% TODO(you): compare to the measured gap (TP=4 step time minus DP=4 step
% time). Does ~7-8 s of exposed TP communication account for it?

% =====================================================================
\section{Part 3: Pipeline Parallel}

\texttt{one\_forward\_one\_backward} runs in three phases. In the warm-up
phase, stage $i$ runs $\min(p-1-i,\,m)$ forward passes. In the steady
state, each iteration runs one forward pass and then the backward pass of
the oldest outstanding micro-batch. This phase uses the combined
\texttt{send\_forward\_recv\_backward} and
\texttt{send\_backward\_recv\_forward} transfers to avoid deadlock. In the
cool-down phase, the remaining backward passes run. A stage never holds
more than $p-i \le p$ micro-batches of activations, whereas AFAB holds
all $m$.

\begin{table}[htbp]
\caption{\texttt{large} at $p=4$, global batch 32, one \texttt{gpuA40x4}
node. Predicted bubble is $(p-1)/(m+p-1)$. Measured idle is
$1-\text{busy\_fraction}$.}
\label{tab:pp}
\centering
\begin{tabular}{rrlrrrr}
\toprule
Micro & $m$ & Schedule & Tokens/s & Predicted bubble & Measured idle & Peak memory (GiB) \\
\midrule
8 & 4  & AFAB & 28{,}802 & 42.9\% & 54.0\% & 17.12 \\
8 & 4  & 1F1B & 28{,}882 & 42.9\% & 54.1\% & 17.12 \\
4 & 8  & AFAB & 32{,}259 & 27.3\% & 43.3\% & 18.42 \\
4 & 8  & 1F1B & 32{,}321 & 27.3\% & 43.1\% & 11.55 \\
2 & 16 & AFAB & 32{,}288 & 15.8\% & 34.8\% & 21.12 \\
2 & 16 & 1F1B & 32{,}246 & 15.8\% & 34.7\% & 8.36  \\
1 & 32 & AFAB & \multicolumn{4}{c}{OOM (44.4\,GiB exhausted on the last stage)} \\
1 & 32 & 1F1B & 27{,}509 & 8.6\%  & 27.9\% & 6.71  \\
\bottomrule
\end{tabular}
\end{table}
% TODO(you): fill from the "bubble and memory" table at the end of
% logs/a4-pp-a40-<id>.out. Write "OOM" in any row that ran out of memory.

\paragraph{Q1: How does the bubble change with $m$?}
% TODO(you): the warm-up and cool-down idle is fixed at (p-1) slots per
% stage. Increasing m adds busy slots but no idle ones, so the fraction
% falls as (p-1)/(m+p-1). Compare the measured idle column with the
% predicted one. If measured > predicted, the extra is P2P transfer time
% and the last stage's heavier work (the output head), which the formula
% ignores.

\paragraph{Q2: What does larger $m$ cost in memory?}
% TODO(you): AFAB keeps all m micro-batches' activations alive until
% backward starts. The last stage also keeps the (micro, T, 50304)
% logits. Their product m x micro = 32 is fixed, so the logits memory
% stays the same while the per-block activations change. 1F1B keeps at
% most p, so its peak should stay roughly flat as m grows. Point at the
% peak-memory column for both schedules.

% =====================================================================
\section{Part 4: Hybrid Parallelism}

\texttt{huge} has 4.04B parameters. With fp32 weights, gradients and two
AdamW moments, that is $16\,\text{B} \times 4.04\times10^9 = 60.1$\,GiB
before any activations, which is more than one 44\,GiB A40 holds. Data
parallelism replicates this state on every rank, so only TP and PP
reduce it. The state per GPU is about $60.1/(tp \cdot pp)$\,GiB:

\begin{center}
\begin{tabular}{rrl}
\toprule
$tp \cdot pp$ & State per GPU & Expectation \\
\midrule
1 & 60.1\,GiB & cannot fit \\
2 & 30.1\,GiB & tight once activations are added \\
4 & 15.0\,GiB & fits \\
8 & 7.5\,GiB  & fits, but no data parallelism is left \\
\bottomrule
\end{tabular}
\end{center}

With the default axis order (\texttt{pp/dp/tp}), TP is the innermost axis,
so a TP group always stays inside one node. That is the right placement
because TP communicates twice per block and nothing hides it. PP is the
outermost axis, so its boundary is the axis that crosses the two nodes.
PP sends only one activation per micro-batch across a stage boundary, so
it puts the least traffic on the slower inter-node link.

\begin{table}[htbp]
\caption{\texttt{huge} on two \texttt{gpuA40x4} nodes (8 A40s), global
batch 128, 6 steps (1 warm-up).}
\label{tab:part4}
\centering
\begin{tabular}{rrrlrrrr}
\toprule
DP & TP & PP & Schedule / DDP & Step time (s) & Tokens/s & Peak memory (GiB) & Loss \\
\midrule
8 & 1 & 1 & --- / naive        & \multicolumn{4}{c}{OOM} \\
2 & 2 & 2 & 1F1B / overlapped  & 40.812 & 3{,}212  & 22.59 & 11.3796 \\
4 & 2 & 1 & --- / overlapped   & 12.811 & 10{,}231 & 38.15 & 11.3773 \\
2 & 1 & 4 & 1F1B / overlapped  & 23.079 & 5{,}679  & 27.58 & 11.3788 \\
4 & 1 & 2 & 1F1B / overlapped  & \tbd & \tbd & \tbd & \tbd \\
2 & 4 & 1 & --- / overlapped   & \tbd & \tbd & \tbd & \tbd \\
1 & 4 & 2 & 1F1B / ---         & \tbd & \tbd & \tbd & \tbd \\
4 & 2 & 1 & --- / naive        & \tbd & \tbd & \tbd & \tbd \\
\bottomrule
\end{tabular}
\end{table}
% TODO(you): fill from out/huge-<tag>/summary.json. The OOM row: paste the
% one line of the CUDA OOM message from logs/a4-part4-<id>.out into the
% text below. Add a row for every other configuration you ran.

\paragraph{Best configuration.}
% TODO(you): name the winner and explain it with the three costs:
%  - what each sharding degree bought in memory (60.1/(tp*pp) + activations),
%  - what it cost in time: TP all-reduces on the critical path inside a node;
%    PP bubble (p-1)/(m+p-1) with m = 128/(dp*micro); DP gradient all-reduce
%    of the state left after TP/PP sharding, overlapped with backward,
%  - why the alternatives you tried lost.
% If you only ran one configuration that fit, say so plainly. Then argue
% from Parts 1-3 which neighbour you would have tried next and why.

\end{document}
